% Round 15. Universal graphic-matroid kernel; independent audit PASS.
\section{A universal graphic-matroid kernel for all split tree types}
\label{r15ug:section}

The fixed affine-tree construction suggests a stronger local device.
One can encode every split marked tree in a single graphic matroid
family, and cut that family out by ambient delta forms before any
input is intersected. This section separates that polyhedral
construction from the still separate ordinary presentation problem.

\begin{definition}[The universal split-tree fans]
\label{r15ug:fans}
For $s\geq2$, let $B_s\subset\mathbf R^{\binom s2}$ be the
weight-one Bergman fan of the graphic matroid of the complete graph
$K_s$, in the minimum valuation convention. Its coordinates are
$\lambda_{ij}$, and the minimum on each cycle occurs at least twice.
For a complete graph it is enough to impose this on triangles.
Let
\[
 C_s=B_{s+1}\subset B_s\times\mathbf R^s,
       \qquad f:C_s\longrightarrow B_s
\tag{UG15.1}\label{r15ug:universal}
\]
forget the edges joining the new vertex to the old ones; denote
those new coordinates by $w_i$. The convention is the negative
of the fan-coordinate convention in Fran\c{c}ois \cite{Francois}.
Every fan, function, and intersection from that reference is
transported by the linear automorphism $u\mapsto-u$ when used here.
For $s=1$ put $B_1=\{0\}$ and $C_1=\mathbf R$.
\end{definition}

\begin{lemma}[The weighted universal tree and flatness]
\label{r15ug:flat}
The fiber of $f$ over $\lambda\in B_s$ is exactly the split
marked tree parametrized by
\[
          w_j=\min(t,\lambda_{ij}),\quad w_i=t,
\tag{UG15.2}\label{r15ug:rays}
\]
with the least-label identification of Lemma~\ref{r15ray:labels}.
All relative edge weights are one, including at degenerations.
The linear map $f$ is flat in Mihatsch's sense.
\end{lemma}

\begin{proof}
The new triangle relations say that the minimum of
$w_i,w_j,\lambda_{ij}$ occurs at least twice. Choose $i$ with
$w_i$ maximal. If $w_i>w_j$, the relation forces
$w_j=\lambda_{ij}$; if $w_i=w_j$, it forces
$w_i\leq\lambda_{ij}$. These are precisely
\eqref{r15ug:rays}. Conversely those values satisfy all new
triangle relations by the ultrametric inequalities for $\lambda$.
The remaining cycles follow from the triangle relations: repeatedly
insert a chord into a cycle having a putative unique minimum and
obtain a shorter cycle with a unique minimum. Thus the description
is the full fiber, not just a selected skeleton inside it.

On every open tree edge, the old coordinates together with the
primitive coordinate $w_i=t$ give the lattice of that cell as the
base-cell lattice times $\mathbf Z$. All other $w_j$ are old
coordinates or $t$. Conversely the inverse coordinate projection
is integral. Thus there is no lattice index, and the quotient
weight is one. This agrees with the weight-one flags-of-flats
calibration of both graphic Bergman fans.

For every compactly supported PWS input on $C_s$, the labelled-band
formula of Proposition~\ref{r15ray:polyhedral} gives a PWS tensor
on $B_s$ with its exact fiber integral against every PWS test.
It includes the tangential restrictions on all faces and the zero
contribution of collapsed intervals. This is exactly Definition
2.18 of Mihatsch \cite{Mihatsch}: existence of PWS fiber integration
for every compact PWS input is the definition of flatness. It also
proves the statement after taking a product with an arbitrary
Euclidean space of extra base coordinates.
\end{proof}

\begin{lemma}[An ambient relative delta kernel]
\label{r15ug:kernel}
There is a global ambient delta-cycle $\Gamma_s$ of bidegree
$(s-1,s-1)$ on
$\mathbf R^{\binom s2}\times\mathbf R^s$ such that
\[
 \Gamma_s\wedge[B_s\times\mathbf R^s]=[C_s].
\tag{UG15.3}\label{r15ug:cut}
\]
More strongly, for every local weighted tropical cycle $T$ with
rational polyhedral directions in
$B_s\times\mathbf R^a$, with its original weights, one has
\[
 \widetilde\Gamma_s\wedge\operatorname{pr}^*[T]
                   =(f\times\mathrm{id})^*[T]
\tag{UG15.4}\label{r15ug:all-cycles}
\]
as polyhedral currents in the corresponding ambient space.
Here $\widetilde\Gamma_s$ ignores the extra $a$ coordinates,
and the right side is the weighted flat pullback, represented
in the universal family. Negative or real cycle weights are
allowed by linearity: on a fixed rational subdivision the balancing
relations have integer coefficients, so their real solution space
is spanned by rational, hence rescaled integer, solutions.
\end{lemma}

\begin{proof}
For $s=1$ take $\Gamma_1=1$. Otherwise put
\[
 M=\operatorname{Mat}(K_s)\oplus U_{s,s},\qquad
 N=\operatorname{Mat}(K_{s+1})
\]
on the common set of old and new graph edges. The new edges
are coloops in $M$. Every circuit of $M$ is an old circuit of
$N$, so $B(N)\subset B(M)=B_s\times\mathbf R^s$.
These are weight-one balanced fans and
$\operatorname{rank}M-\operatorname{rank}N=s-1$.
Fran\c{c}ois \cite{Francois}, Remark 4.4, gives piecewise linear
functions $\phi_1,\ldots,\phi_{s-1}$ on $B(M)$ whose divisor
product is $B(N)$. In that paper's sign convention their value
at the flat ray $V_F$ is $-1$ precisely when
$r_M(F)-r_N(F)\geq i$, and zero otherwise. Transport the formula
to our convention.

Extend each of these finite-fan piecewise linear functions to
an ambient continuous piecewise linear function $\Phi_i$.
Such an extension is obtained by extending the finite fan to a
complete subdivision and assigning values on its added rays,
while extending the common linear part on the lineality space.
Set
\[
                       \Gamma_s=D\Phi_1\wedge\cdots
                                            \wedge D\Phi_{s-1}.
\tag{UG15.5}\label{r15ug:explicit}
\]
The Euclidean delta intersection agrees with successive tropical
Cartier divisors, by Mihatsch \cite{MihatschIntersection},
Theorem 4.1. Consequently \eqref{r15ug:cut} is an identity of
weighted cycles with the stated normalization.

For \eqref{r15ug:all-cycles}, work locally near any point of
$B_s\times\mathbf R^a$. Its star is a product of matroid-minor
fans and a linear space. The local fan form of Fran\c{c}ois,
Theorem 4.1 and Remark 4.2, cuts the local cycle $T$ out by a
piecewise polynomial cocycle on this smooth tropical space.
Piecewise polynomials on a finite fan are sums of products of
piecewise linear functions (loc. cit., Section 3). Extending
these factors to the Euclidean ambient space gives an ambient
delta-form $\zeta$ such that
\[
                 [T]=\zeta\wedge[B_s\times\mathbf R^a]
\tag{UG15.6}\label{r15ug:local-cut}
\]
on the chosen local open. This uses only local stars; no assertion
about gluing global cocycles of arbitrary degree is needed.
In particular $T$, intrinsically on that smooth tropical space,
is a Mihatsch delta-form.

Associativity and \eqref{r15ug:cut} now give
\[
 \widetilde\Gamma_s\wedge\operatorname{pr}^*[T]
       =\operatorname{pr}^*\zeta\wedge[C_s\times\mathbf R^a].
\tag{UG15.7}\label{r15ug:commute}
\]
By Lemma~\ref{r15ug:flat} and Mihatsch Proposition 2.44, flat
pullback of the intrinsic delta-form $T$ agrees with substitution
of its genuine ambient presentation. The right side of
\eqref{r15ug:commute} is exactly that pullback. This proves
\eqref{r15ug:all-cycles}. Equivalently one can apply the
successive divisor formulas to the ambient PL factors of $\zeta$.
No general delta pullback along the nonflat inclusion of the
universal family is asserted. Local equality of these currents
glues by uniqueness, without a global cocycle-injectivity claim.
\end{proof}

\begin{theorem}[Universal output for a supplied split ambient presentation]
\label{r15ug:projection}
Keep the smooth algebraic base, split analytic root functions,
and complete tuple of difference units of Definition~\ref{r15ft:data},
but impose no fixed tree-type condition. Suppose the input still
has its supplied single ambient presentation with proper ambient
vertical support. Then its projective-line current has an ordinary
actual representative of degree $(p-1,q-1)$, with every ordinary
smooth compact mixed-test pairing and the native normalization.
\end{theorem}

\begin{proof}
Separate the root-difference coordinates from the extra base
coordinates, duplicating a coordinate if necessary by a genuine
unit refinement. Pull the ambient kernel $\Gamma_s$ back by that
coordinate projection. For the supplied ambient input $\omega$ put
\[
          \beta_0=\operatorname{pr}_*(\Gamma_s\wedge\omega),
                   \qquad\beta=t_b^\star\beta_0.
\tag{UG15.8}\label{r15ug:output}
\]
Proper ambient support permits Euclidean affine delta pushforward;
its degree is $(p+s-1-s,q+s-1-s)=(p-1,q-1)$.
Thus $\beta$ is an ordinary actual form on the base.

Append every target-test unit and the compact-localization cutoff
moments, as in Theorem~\ref{r15ft:projection}. The refined base
fundamental cycle $T$ lies in $B_s$ times the Euclidean space of
all extra coordinates. On compact algebraic Laurent subcharts,
exact-norm rationalization of the finite base tuple makes these
fundamental cycles rational in their directions; arbitrary real
domain radii do not change that assertion. The maximal initial-form argument in that
theorem identifies the source fundamental cycle with the marked-tree
family over $T$, with base weight times relative weight one.
Lemma~\ref{r15ug:kernel} identifies that very whole cycle with
$\Gamma_s\wedge\operatorname{pr}^*T$. This equality precedes
intersection with the input, so retains every positive-delta stratum
at vertices and changing tree types. The Euclidean projection formula
and Cartesian product by the extra test coordinates prove the
prescribed pairing with \eqref{r15ug:output}. Ordinary current
injectivity proves uniqueness and gluing. Finite descent is valid
whenever the output is in the established actual finite-trace strip,
with division by the full native degree.
\end{proof}

\paragraph{The separate analytic presentation interface.}
The theorem requires a supplied single ambient actual presentation.
The earlier PWS assembly by degreewise PWS extension cannot simply
be reused for positive-delta inputs. Proposition~\ref{r15ad:presenter}
supplies the necessary new construction using two nested Laurent
neighborhoods, strict numerical margins, and genuine Euclidean
cutoffs. Combining that proposition with this kernel and the frozen
finite actual trace gives Theorem~\ref{r15a2:projection} for
$\min(p,q)\leq2$. The universal kernel itself works in all degrees
for supplied ambient inputs, but descent of a general output beyond
the established finite-trace strip remains open.
